\sectionheading{5.40 Python \textenglish{\ttfamily while} Loop}

কোনো Condition যতক্ষণ
\textenglish{\ttfamily True}
থাকে, ততক্ষণ একই Code Block বারবার Execute করার জন্য Python-এ
\textenglish{\ttfamily while}
Loop ব্যবহার করা হয়।

%================================================
% Definition
%================================================

\booksubheading{\textenglish{\ttfamily while} Loop কী?}

\begin{definitionbox}

কোনো নির্দিষ্ট Condition
\textenglish{\ttfamily True}
থাকা পর্যন্ত একটি Statement বা Code Block পুনরাবৃত্তি করার Loop-কে
\textbf{\textenglish{\ttfamily while} Loop} বলে।

\end{definitionbox}

\begin{center}

\textbf{Condition পরীক্ষা}
\(\longrightarrow\)
\textenglish{\ttfamily True}
\(\longrightarrow\)
\textbf{Loop Body}
\(\longrightarrow\)
\textbf{Update}

\medskip

\textbf{Update}
\(\longrightarrow\)
\textbf{আবার Condition পরীক্ষা}

\end{center}

Condition
\textenglish{\ttfamily False}
হলে Loop শেষ হয়।

%================================================
% Syntax
%================================================

\booksubheading{\textenglish{\ttfamily while} Loop-এর Syntax}

\begin{codebox}
\begin{english}
\ttfamily
initialization\\[0.5mm]
while condition:\\
\hspace*{4mm}statement\\
\hspace*{4mm}update
\end{english}
\end{codebox}

উদাহরণ—

\begin{codebox}
\begin{english}
\ttfamily
i = 1\\[0.5mm]
while i <= 5:\\
\hspace*{4mm}print(i)\\
\hspace*{4mm}i += 1
\end{english}
\end{codebox}

\textbf{Output:}

\begin{center}
\textenglish{\ttfamily 1 \quad 2 \quad 3 \quad 4 \quad 5}
\end{center}

%================================================
% Main parts
%================================================

\booksubheading{\textenglish{\ttfamily while} Loop-এর প্রধান অংশ}

\begin{tabularx}
{\linewidth}
{
|>{\bfseries\RaggedRight\arraybackslash}p{0.25\linewidth}
|>{\RaggedRight\arraybackslash}p{0.27\linewidth}
|Y|
}

\hline
\rowcolor{softBlue}

\textbf{অংশ}
&
\textbf{Code}
&
\textbf{কাজ}
\\
\hline

Initialization
&
\textenglish{\ttfamily i = 1}
&
প্রাথমিক Value নির্ধারণ করে।
\\
\hline

Condition
&
\textenglish{\ttfamily i <= 5}
&
Loop চলবে কি না পরীক্ষা করে।
\\
\hline

Loop Body
&
\textenglish{\ttfamily print(i)}
&
পুনরাবৃত্ত কাজ সম্পন্ন করে।
\\
\hline

Update
&
\textenglish{\ttfamily i += 1}
&
Counter-এর Value পরিবর্তন করে।
\\
\hline

\end{tabularx}

\begin{rememberbox}

\textbf{\color{deepOrange}মনে রাখবে}

\begin{center}

\textbf{Initialization}
\(\longrightarrow\)
\textbf{Condition}
\(\longrightarrow\)
\textbf{Loop Body}
\(\longrightarrow\)
\textbf{Update}

\end{center}

এ চারটি অংশ
\textenglish{\ttfamily while}
Loop বোঝার জন্য সবচেয়ে গুরুত্বপূর্ণ।

\end{rememberbox}

%================================================
% Working process
%================================================

\booksubheading{\textenglish{\ttfamily while} Loop কীভাবে কাজ করে?}

ধরা যাক—

\begin{codebox}
\begin{english}
\ttfamily
i = 1\\[0.5mm]
while i <= 3:\\
\hspace*{4mm}print(i)\\
\hspace*{4mm}i += 1
\end{english}
\end{codebox}

এর কার্যধারা—

\begin{center}

\textenglish{\ttfamily i = 1}
\(\longrightarrow\)
\textenglish{\ttfamily 1 <= 3}
True
\(\longrightarrow\)
Print 1
\(\longrightarrow\)
\textenglish{\ttfamily i = 2}

\medskip

\textenglish{\ttfamily i = 2}
\(\longrightarrow\)
\textenglish{\ttfamily 2 <= 3}
True
\(\longrightarrow\)
Print 2
\(\longrightarrow\)
\textenglish{\ttfamily i = 3}

\medskip

\textenglish{\ttfamily i = 3}
\(\longrightarrow\)
\textenglish{\ttfamily 3 <= 3}
True
\(\longrightarrow\)
Print 3
\(\longrightarrow\)
\textenglish{\ttfamily i = 4}

\medskip

\textenglish{\ttfamily i = 4}
\(\longrightarrow\)
\textenglish{\ttfamily 4 <= 3}
False
\(\longrightarrow\)
\textbf{Loop শেষ}

\end{center}

%================================================
% Trace table
%================================================

\booksubheading{Dry Run / Trace Table}

উপরের Program-এর Trace Table—

\begin{tabularx}
{\linewidth}
{
|>{\bfseries\centering\arraybackslash}p{0.16\linewidth}
|>{\centering\arraybackslash}p{0.13\linewidth}
|>{\centering\arraybackslash}p{0.21\linewidth}
|>{\centering\arraybackslash}p{0.17\linewidth}
|>{\centering\arraybackslash}Y|
}

\hline
\rowcolor{softBlue}

\textbf{Iteration}
&
\textbf{\textenglish{\ttfamily i}}
&
\textbf{\textenglish{\ttfamily i <= 3}}
&
\textbf{Output}
&
\textbf{Update-এর পর \textenglish{\ttfamily i}}
\\
\hline

1
&
1
&
True
&
1
&
2
\\
\hline

2
&
2
&
True
&
2
&
3
\\
\hline

3
&
3
&
True
&
3
&
4
\\
\hline

শেষ পরীক্ষা
&
4
&
False
&
—
&
—
\\
\hline

\end{tabularx}

\begin{rememberbox}

\textbf{\color{deepOrange}MCQ তথ্য:}
Condition
\textenglish{\ttfamily False}
হওয়ার পর Loop Body আর Execute হয় না।

\end{rememberbox}

%================================================
% Forward Reverse
%================================================

\booksubheading{Forward ও Reverse \textenglish{\ttfamily while} Loop}

\noindent
\begin{minipage}[t]{0.49\linewidth}

\begin{sideexamplebox}{Forward Loop}

ছোট Value থেকে বড় Value-এর দিকে যেতে Counter বাড়ানো হয়।

\begin{sidecodebox}
\begin{english}
\ttfamily
i = 1\\
while i <= 5:\\
\hspace*{4mm}print(i)\\
\hspace*{4mm}i += 1
\end{english}
\end{sidecodebox}

\textbf{Output:}

\begin{center}
\textenglish{\ttfamily 1 \quad 2 \quad 3 \quad 4 \quad 5}
\end{center}

\end{sideexamplebox}

\end{minipage}
\hfill
\begin{minipage}[t]{0.49\linewidth}

\begin{sideexamplebox}{Reverse Loop}

বড় Value থেকে ছোট Value-এর দিকে যেতে Counter কমানো হয়।

\begin{sidecodebox}
\begin{english}
\ttfamily
i = 5\\
while i >= 1:\\
\hspace*{4mm}print(i)\\
\hspace*{4mm}i -= 1
\end{english}
\end{sidecodebox}

\textbf{Output:}

\begin{center}
\textenglish{\ttfamily 5 \quad 4 \quad 3 \quad 2 \quad 1}
\end{center}

\end{sideexamplebox}

\end{minipage}

\begin{rememberbox}

\begin{center}

\textbf{Forward}
\(\longrightarrow\)
সাধারণত
\textenglish{\ttfamily i += 1}

\qquad

\textbf{Reverse}
\(\longrightarrow\)
সাধারণত
\textenglish{\ttfamily i -= 1}

\end{center}

\end{rememberbox}

%================================================
% Even Odd
%================================================

\booksubheading{Even ও Odd Number}

\noindent
\begin{minipage}[t]{0.49\linewidth}

\begin{sideexamplebox}{Even Number}

২ থেকে ১০ পর্যন্ত জোড় সংখ্যা—

\begin{sidecodebox}
\begin{english}
\ttfamily
i = 2\\
while i <= 10:\\
\hspace*{4mm}print(i)\\
\hspace*{4mm}i += 2
\end{english}
\end{sidecodebox}

Output:

\begin{center}
\textenglish{\ttfamily 2, 4, 6, 8, 10}
\end{center}

\end{sideexamplebox}

\end{minipage}
\hfill
\begin{minipage}[t]{0.49\linewidth}

\begin{sideexamplebox}{Odd Number}

১ থেকে ১০ পর্যন্ত বিজোড় সংখ্যা—

\begin{sidecodebox}
\begin{english}
\ttfamily
i = 1\\
while i <= 10:\\
\hspace*{4mm}print(i)\\
\hspace*{4mm}i += 2
\end{english}
\end{sidecodebox}

Output:

\begin{center}
\textenglish{\ttfamily 1, 3, 5, 7, 9}
\end{center}

\end{sideexamplebox}

\end{minipage}

\begin{examplebox}{Condition ব্যবহার করেও Even Number বের করা যায়}

\begin{codebox}
\begin{english}
\ttfamily
i = 1\\
while i <= 10:\\
\hspace*{4mm}if i \% 2 == 0:\\
\hspace*{8mm}print(i)\\
\hspace*{4mm}i += 1
\end{english}
\end{codebox}

\end{examplebox}

%================================================
% Sum
%================================================

\booksubheading{Loop ব্যবহার করে Sum নির্ণয়}

১ থেকে ৫ পর্যন্ত সংখ্যার যোগফল—

\begin{codebox}
\begin{english}
\ttfamily
i = 1\\
total = 0\\[0.5mm]
while i <= 5:\\
\hspace*{4mm}total += i\\
\hspace*{4mm}i += 1\\[0.5mm]
print(total)
\end{english}
\end{codebox}

\textbf{Output:}
\textenglish{\ttfamily 15}

\begin{center}

\textenglish{\ttfamily 0}
\(\longrightarrow\)
\textenglish{\ttfamily +1}
\(\longrightarrow\)
\textenglish{\ttfamily +2}
\(\longrightarrow\)
\textenglish{\ttfamily +3}
\(\longrightarrow\)
\textenglish{\ttfamily +4}
\(\longrightarrow\)
\textenglish{\ttfamily +5}
\(\longrightarrow\)
\textenglish{\ttfamily 15}

\end{center}

\begin{tabularx}
{\linewidth}
{
|>{\bfseries\RaggedRight\arraybackslash}p{0.32\linewidth}
|Y|
}

\hline
\rowcolor{softBlue}

\textbf{Code}
&
\textbf{কাজ}
\\
\hline

\textenglish{\ttfamily i = 1}
&
Counter-এর Initial Value।
\\
\hline

\textenglish{\ttfamily total = 0}
&
Sum-এর Initial Value।
\\
\hline

\textenglish{\ttfamily total += i}
&
প্রতিটি Iteration-এ Sum Update।
\\
\hline

\textenglish{\ttfamily i += 1}
&
Counter Update।
\\
\hline

\end{tabularx}

%================================================
% Fixed number repetition
%================================================

\booksubheading{নির্দিষ্ট সংখ্যকবার Statement Execute}

\begin{examplebox}{\textenglish{\ttfamily "Python"} পাঁচবার Print}

\begin{codebox}
\begin{english}
\ttfamily
count = 1\\[0.5mm]
while count <= 5:\\
\hspace*{4mm}print("Python")\\
\hspace*{4mm}count += 1
\end{english}
\end{codebox}

এখানে Loop মোট \textbf{5 Iteration} চলবে।

\end{examplebox}

%================================================
% Condition based input
%================================================

\booksubheading{Condition-ভিত্তিক Loop}

\textenglish{\ttfamily while}
Loop-এর বড় সুবিধা হলো—কতবার Loop চলবে তা আগে থেকে জানা না থাকলেও
Condition-এর ভিত্তিতে Loop চালানো যায়।

\begin{examplebox}{User Input-এর উপর Loop}

\begin{codebox}
\begin{english}
\ttfamily
number = int(input("Enter 0 to stop: "))\\[0.5mm]
while number != 0:\\
\hspace*{4mm}print("You entered:", number)\\
\hspace*{4mm}number = int(input("Enter 0 to stop: "))
\end{english}
\end{codebox}

User যতক্ষণ
\textenglish{\ttfamily 0}
না দেয়, Loop চলবে।

\end{examplebox}

%================================================
% Sentinel
%================================================

\booksubheading{Sentinel Value}

Loop বন্ধ করার জন্য বিশেষভাবে ব্যবহৃত Value-কে অনেক ক্ষেত্রে
\textbf{Sentinel Value} বলা হয়।

উপরের Program-এ—

\begin{center}

\textbf{Sentinel Value}
\(\longrightarrow\)
\textenglish{\ttfamily 0}

\end{center}

\begin{rememberbox}

Sentinel-controlled Loop-এ বিশেষ একটি Value পাওয়া গেলে
Loop বন্ধ করা হয়।

\end{rememberbox}

%================================================
% Update importance
%================================================

\booksubheading{\textenglish{\ttfamily while} Loop-এ Update কেন গুরুত্বপূর্ণ?}

\noindent
\begin{minipage}[t]{0.49\linewidth}

\begin{sideexamplebox}{সঠিক}

\begin{sidecodebox}
\begin{english}
\ttfamily
i = 1\\
while i <= 5:\\
\hspace*{4mm}print(i)\\
\hspace*{4mm}i += 1
\end{english}
\end{sidecodebox}

Counter পরিবর্তিত হচ্ছে, তাই Condition একসময় False হবে।

\end{sideexamplebox}

\end{minipage}
\hfill
\begin{minipage}[t]{0.49\linewidth}

\begin{sideexamplebox}{ভুল}

\begin{sidecodebox}
\begin{english}
\ttfamily
i = 1\\
while i <= 5:\\
\hspace*{4mm}print(i)
\end{english}
\end{sidecodebox}

\textenglish{\ttfamily i}
পরিবর্তিত হচ্ছে না, তাই Loop চলতেই থাকবে।

\end{sideexamplebox}

\end{minipage}

\begin{rememberbox}

\textbf{\color{deepOrange}ভুল করো না}

\begin{center}

\textbf{Counter Update না করলে}
\(\longrightarrow\)
\textbf{Infinite Loop হতে পারে}

\end{center}

\end{rememberbox}

%================================================
% Infinite Loop
%================================================

\booksubheading{Infinite Loop}

যে Loop-এর Condition কার্যকরভাবে কখনো
\textenglish{\ttfamily False}
হয় না এবং Loop অনির্দিষ্টকাল চলতে থাকে তাকে
\textbf{Infinite Loop} বলে।

\begin{codebox}
\begin{english}
\ttfamily
while True:\\
\hspace*{4mm}print("Python")
\end{english}
\end{codebox}

এটি নিজে থেকে শেষ হবে না।

\begin{tabularx}
{\linewidth}
{
|>{\bfseries\RaggedRight\arraybackslash}p{0.36\linewidth}
|Y|
}

\hline
\rowcolor{softBlue}

\textbf{সাধারণ কারণ}
&
\textbf{উদাহরণ}
\\
\hline

Counter Update না করা
&
\textenglish{\ttfamily i}
-এর Value অপরিবর্তিত থাকা
\\
\hline

ভুল Direction-এ Update
&
কমার কথা থাকলে বাড়ানো
\\
\hline

Condition কখনো False না হওয়া
&
\textenglish{\ttfamily while True}
\\
\hline

\end{tabularx}

%================================================
% Intentional infinite loop
%================================================

\begin{rememberbox}

\textbf{\color{deepOrange}সব Infinite Loop ভুল নয়}

কখনো ইচ্ছাকৃতভাবে—

\begin{center}
\textenglish{\ttfamily while True:}
\end{center}

ব্যবহার করা হয়। প্রয়োজনমতো
\textenglish{\ttfamily break}
দিয়ে Loop বন্ধ করা যায়।

\textenglish{\ttfamily break}
বিস্তারিতভাবে
\textbf{5.41 Loop Control Statement}
-এ আলোচনা হবে।

\end{rememberbox}

%================================================
% Zero iteration + entry controlled
%================================================

\booksubheading{Condition শুরুতেই False হলে কী হবে?}

\begin{codebox}
\begin{english}
\ttfamily
i = 10\\[0.5mm]
while i < 5:\\
\hspace*{4mm}print(i)
\end{english}
\end{codebox}

এখানে শুরুতেই—

\begin{center}

\textenglish{\ttfamily 10 < 5}
\(\longrightarrow\)
\textenglish{\ttfamily False}

\end{center}

তাই Loop Body
\textbf{একবারও Execute হবে না}।

\begin{rememberbox}

\textbf{\color{deepOrange}MCQ Trap}

\textenglish{\ttfamily while}
হলো \textbf{Entry-Controlled Loop}।

কারণ Loop Body-তে প্রবেশের আগেই Condition পরীক্ষা করা হয়।

\begin{center}

\textbf{শুরুতেই Condition False}
\(\longrightarrow\)
\textbf{0 Iteration}

\end{center}

\end{rememberbox}

%================================================
% Multiple statements and indentation
%================================================

\booksubheading{Loop Body ও Indentation}

\noindent
\begin{minipage}[t]{0.49\linewidth}

\begin{sideexamplebox}{একাধিক Statement Loop Body-তে}

\begin{sidecodebox}
\begin{english}
\ttfamily
i = 1\\
while i <= 3:\\
\hspace*{4mm}print("Number:", i)\\
\hspace*{4mm}print("ICT")\\
\hspace*{4mm}i += 1
\end{english}
\end{sidecodebox}

Indented তিনটি Statement-ই Loop Body-এর অংশ।

\end{sideexamplebox}

\end{minipage}
\hfill
\begin{minipage}[t]{0.49\linewidth}

\begin{sideexamplebox}{Statement Loop-এর বাইরে}

\begin{sidecodebox}
\begin{english}
\ttfamily
i = 1\\
while i <= 3:\\
\hspace*{4mm}print(i)\\
\hspace*{4mm}i += 1\\[0.5mm]
print("Finished")
\end{english}
\end{sidecodebox}

শেষের Statement মাত্র একবার Execute হবে।

\end{sideexamplebox}

\end{minipage}

\begin{rememberbox}

Python-এ
\textenglish{\ttfamily \{ \}}
দিয়ে Loop Body নির্ধারণ করা হয় না;
\textbf{Indentation} দিয়ে Code Block নির্ধারিত হয়।

\end{rememberbox}

%================================================
% Update variations
%================================================

\booksubheading{Update সবসময় \textenglish{\ttfamily +1} হতে হবে?}

না। সমস্যার Logic অনুযায়ী Update ভিন্ন হতে পারে।

\begin{tabularx}
{\linewidth}
{
|>{\bfseries\RaggedRight\arraybackslash}p{0.30\linewidth}
|Y|
}

\hline
\rowcolor{softBlue}

\textbf{Update}
&
\textbf{অর্থ}
\\
\hline

\textenglish{\ttfamily i += 1}
&
১ করে বৃদ্ধি।
\\
\hline

\textenglish{\ttfamily i += 2}
&
২ করে বৃদ্ধি।
\\
\hline

\textenglish{\ttfamily i -= 1}
&
১ করে হ্রাস।
\\
\hline

\textenglish{\ttfamily i *= 2}
&
প্রতিবার ২ দিয়ে গুণ।
\\
\hline

\end{tabularx}

Condition এবং Problem-এর Logic অনুযায়ী সঠিক Update নির্বাচন করতে হয়।

%================================================
% Counter
%================================================

\booksubheading{\textenglish{\ttfamily while} Loop-এ Counter}

\begin{codebox}
\begin{english}
\ttfamily
count = 1\\[0.5mm]
while count <= 3:\\
\hspace*{4mm}print(count)\\
\hspace*{4mm}count += 1
\end{english}
\end{codebox}

এখানে
\textenglish{\ttfamily count}
হলো Counter।

Counter সাধারণত—

\begin{center}

\textbf{Loop-এর আগে Initialize}
\(\longrightarrow\)
\textbf{Condition-এ ব্যবহার}
\(\longrightarrow\)
\textbf{Loop Body-তে Update}

\end{center}

%================================================
% Nested while
%================================================

\booksubheading{Nested \textenglish{\ttfamily while} Loop}

একটি
\textenglish{\ttfamily while}
Loop-এর ভিতরে আরেকটি
\textenglish{\ttfamily while}
Loop থাকলে তাকে
\textbf{Nested while Loop} বলে।

\begin{examplebox}{উদাহরণ}

\begin{codebox}
\begin{english}
\ttfamily
i = 1\\[0.5mm]
while i <= 2:\\
\hspace*{4mm}j = 1\\[0.5mm]
\hspace*{4mm}while j <= 3:\\
\hspace*{8mm}print(i, j)\\
\hspace*{8mm}j += 1\\[0.5mm]
\hspace*{4mm}i += 1
\end{english}
\end{codebox}

এখানে—

\begin{center}

\textbf{Outer Loop}
\(\longrightarrow\)
2 বার

\qquad

\textbf{Inner Loop}
\(\longrightarrow\)
প্রতিবার 3 বার

\end{center}

তাই
\textenglish{\ttfamily print()}
মোট—

\begin{center}

\(2 \times 3 = 6\) বার Execute হবে।

\end{center}

\end{examplebox}

%================================================
% Inner counter note
%================================================

\begin{rememberbox}

\textbf{\color{deepOrange}Nested Loop-এ গুরুত্বপূর্ণ}

Inner Loop-এর Counter সাধারণত Outer Loop-এর প্রতিটি Iteration-এ
আবার Initialize করতে হয়।

\begin{center}

\textenglish{\ttfamily while i <= 2:}
\(\longrightarrow\)
\textenglish{\ttfamily j = 1}

\end{center}

\textenglish{\ttfamily j}
পুনরায় Initialize না করলে পরবর্তী Outer Iteration-এ Inner Loop
প্রয়োজনমতো নাও চলতে পারে।

\end{rememberbox}

%================================================
% while else
%================================================

\booksubheading{\textenglish{\ttfamily while ... else}}

Python-এ
\textenglish{\ttfamily while}
Loop-এর সঙ্গে Optional
\textenglish{\ttfamily else}
ব্যবহার করা যায়।

\begin{codebox}
\begin{english}
\ttfamily
i = 1\\[0.5mm]
while i <= 3:\\
\hspace*{4mm}print(i)\\
\hspace*{4mm}i += 1\\
else:\\
\hspace*{4mm}print("Loop Finished")
\end{english}
\end{codebox}

\textbf{Output:}

\begin{codebox}
\begin{english}
\ttfamily
1\\
2\\
3\\
Loop Finished
\end{english}
\end{codebox}

Loop স্বাভাবিকভাবে Condition
\textenglish{\ttfamily False}
হয়ে শেষ হলে
\textenglish{\ttfamily else}
Block Execute হয়।

\begin{rememberbox}

\textenglish{\ttfamily break}
দিয়ে Loop বন্ধ হলে
\textenglish{\ttfamily while ... else}
-এর
\textenglish{\ttfamily else}
Block Execute হয় না।

এটি পরবর্তী
\textbf{5.41 Loop Control Statement}
Topic-এ আরও পরিষ্কার হবে।

\end{rememberbox}

%================================================
% for vs while
%================================================

\booksubheading{\textenglish{\ttfamily for} ও \textenglish{\ttfamily while}-এর ব্যবহার}

একই কাজ অনেক সময় দুই ধরনের Loop দিয়েই করা যায়।

\noindent
\begin{minipage}[t]{0.49\linewidth}

\begin{sideexamplebox}{\textenglish{\ttfamily for}}

\begin{sidecodebox}
\begin{english}
\ttfamily
for i in range(1, 6):\\
\hspace*{4mm}print(i)
\end{english}
\end{sidecodebox}

নির্দিষ্ট Range-এর জন্য সুবিধাজনক।

\end{sideexamplebox}

\end{minipage}
\hfill
\begin{minipage}[t]{0.49\linewidth}

\begin{sideexamplebox}{\textenglish{\ttfamily while}}

\begin{sidecodebox}
\begin{english}
\ttfamily
i = 1\\
while i <= 5:\\
\hspace*{4mm}print(i)\\
\hspace*{4mm}i += 1
\end{english}
\end{sidecodebox}

Condition-ভিত্তিক Repetition-এর জন্য সুবিধাজনক।

\end{sideexamplebox}

\end{minipage}

\begin{tabularx}
{\linewidth}
{
|>{\bfseries\RaggedRight\arraybackslash}p{0.39\linewidth}
|>{\centering\arraybackslash}Y|
}

\hline
\rowcolor{softBlue}

\textbf{পরিস্থিতি}
&
\textbf{সাধারণত সুবিধাজনক}
\\
\hline

নির্দিষ্ট Range
&
\textenglish{\ttfamily for}
\\
\hline

Sequence/Collection-এর Item Process
&
\textenglish{\ttfamily for}
\\
\hline

Condition True থাকা পর্যন্ত
&
\textenglish{\ttfamily while}
\\
\hline

Iteration সংখ্যা আগে থেকে স্পষ্ট নয়
&
\textenglish{\ttfamily while}
\\
\hline

User Input-এর উপর Repetition
&
\textenglish{\ttfamily while}
\\
\hline

\end{tabularx}

%================================================
% Use cases
%================================================

\booksubheading{\textenglish{\ttfamily while} Loop কোথায় বেশি উপযোগী?}

\begin{compacttwocolpoints}

\item Condition পূরণ হওয়া পর্যন্ত কাজ চালাতে।

\item User-এর Input-এর উপর Loop চালাতে।

\item Valid Input না পাওয়া পর্যন্ত পুনরায় Input নিতে।

\item Menu-driven Program তৈরি করতে।

\item কোনো Target পাওয়া পর্যন্ত Searching করতে।

\item নির্দিষ্ট Event ঘটার অপেক্ষায় Repetition চালাতে।

\item Iteration সংখ্যা আগে থেকে জানা না থাকলে।

\end{compacttwocolpoints}

%================================================
% Input validation
%================================================

\booksubheading{Input Validation-এর সহজ উদাহরণ}

User Positive Number না দেওয়া পর্যন্ত Input নেওয়া—

\begin{examplebox}{Valid Input না পাওয়া পর্যন্ত Loop}

\begin{codebox}
\begin{english}
\ttfamily
n = int(input("Enter positive number: "))\\[0.5mm]
while n <= 0:\\
\hspace*{4mm}n = int(input("Try again: "))\\[0.5mm]
print("Accepted")
\end{english}
\end{codebox}

এখানে Loop কতবার চলবে তা আগে থেকে জানা নেই।
তাই
\textenglish{\ttfamily while}
Loop উপযোগী।

\end{examplebox}

%================================================
% Key features
%================================================

\booksubheading{\textenglish{\ttfamily while} Loop-এর গুরুত্বপূর্ণ বৈশিষ্ট্য}

\begin{tabularx}
{\linewidth}
{
|>{\bfseries\RaggedRight\arraybackslash}p{0.36\linewidth}
|Y|
}

\hline
\rowcolor{softBlue}

\textbf{বিষয়}
&
\textbf{তথ্য}
\\
\hline

Keyword
&
\textenglish{\ttfamily while}
\\
\hline

মূল ভিত্তি
&
Condition
\\
\hline

Condition True
&
Loop Body Execute
\\
\hline

Condition False
&
Loop শেষ
\\
\hline

Code Block
&
Indentation দ্বারা নির্ধারিত
\\
\hline

Initialization
&
সাধারণত Loop-এর আগে
\\
\hline

Update
&
Programmer-কে লিখতে হয়
\\
\hline

0 Iteration সম্ভব?
&
হ্যাঁ
\\
\hline

Infinite Loop সম্ভব?
&
হ্যাঁ
\\
\hline

Loop Control
&
Entry-Controlled
\\
\hline

\textenglish{\ttfamily else}
&
Optional
\\
\hline

\end{tabularx}

%================================================
% Common mistakes
%================================================

\booksubheading{সাধারণ ভুল}

\begin{rememberbox}

\begin{itemize}

\item Loop-এর আগে Counter Initialize না করা।

\item Condition ভুল লেখা।

\item Counter/Variable Update না করা।

\item সঠিক Direction-এর বিপরীতে Update করা।

\item Condition-এর শেষে
\textenglish{\ttfamily :}
না দেওয়া।

\item Loop Body-তে সঠিক Indentation না দেওয়া।

\item ভুল করে Infinite Loop তৈরি করা।

\item Reverse Loop-এ Counter কমানোর পরিবর্তে বাড়ানো।

\item Loop-এর বাইরে থাকা Statement-কে Loop Body মনে করা।

\item Nested
\textenglish{\ttfamily while}
-এ Inner Counter পুনরায় Initialize না করা।

\item শুরুতেই Condition False হলে Loop Body চলবে মনে করা।

\end{itemize}

\end{rememberbox}

%================================================
% Final Reminder
%================================================

\begin{rememberbox}

\textbf{\color{deepOrange}মনে রাখবে}

\begin{tabularx}
{\linewidth}
{
>{\bfseries\RaggedRight\arraybackslash}p{0.24\linewidth}
Y
>{\bfseries\RaggedRight\arraybackslash}p{0.24\linewidth}
Y
}

\textenglish{\ttfamily while}
&
Condition True থাকা পর্যন্ত
&
Initialization
&
Loop-এর আগে
\\[0.7mm]

Condition
&
Loop চলবে কি না
&
Loop Body
&
Repeated Code
\\[0.7mm]

Update
&
Counter পরিবর্তন
&
Condition False
&
Loop শেষ
\\[0.7mm]

Forward
&
Counter বৃদ্ধি
&
Reverse
&
Counter হ্রাস
\\[0.7mm]

0 Iteration
&
সম্ভব
&
Infinite Loop
&
সম্ভব
\\[0.7mm]

Entry-Controlled
&
আগে Condition পরীক্ষা
&
Nested Loop
&
Loop-এর ভিতরে Loop
\\[0.7mm]

Sentinel
&
Loop বন্ধের বিশেষ Value
&
\textenglish{\ttfamily else}
&
Optional
\\

\end{tabularx}

\medskip

\begin{center}

\textbf{Initialization}
\(\longrightarrow\)
\textbf{Condition}

\medskip

\textbf{Condition True}
\(\longrightarrow\)
\textbf{Loop Body}
\(\longrightarrow\)
\textbf{Update}
\(\longrightarrow\)
\textbf{আবার Condition}

\medskip

\textbf{Condition False}
\(\longrightarrow\)
\textbf{Loop শেষ}

\medskip

\textbf{Update না হলে}
\(\longrightarrow\)
\textbf{Infinite Loop হতে পারে}

\end{center}

\end{rememberbox}

\small

%================================================
% 5.41 Loop Control Statement
%================================================

